\chapter{Reading a World Back}
\label{ch:reading}

\section{Why the Reads Are the Interface}

An agent cannot look at the map. Everything it learns about what it has built comes
through a read, and the quality of those reads therefore bounds the quality of the
authoring. This is the chapter that makes the rest of the system operable.

The studio offers about thirty of them. They divide cleanly into pictures and
numbers, and most have both forms: a PNG for a person, and a \id{?format=text}
twin that answers the same question as characters or a table. The rule the corpus
settled on is that the text comes first. A picture answers \emph{whether} something
came out. A number answers \emph{whether it is right}.

\section{Every Read States What It Cannot Answer}

The property that distinguishes this surface from an ordinary set of debugging
endpoints is that each one documents its own blind spot. The studio's published
schema gives every read a summary and, beside it, a clause headed \textsc{misleads}
naming the wrong conclusion a reader is likely to draw from it.

This is unusual enough to deserve examples, because the clauses are not
boilerplate. They are the specific mistakes that have been made.

\begin{quote}
\textbf{Column.} A column through the middle of a house reads floor, air, roof. The
walls are at the perimeter; that is a correct building, not a broken one.
\end{quote}

\begin{quote}
\textbf{Top-down render.} It keeps no Y at all, so a riser, a ramp's step heights, a
stamped room's floor and a goal's clearance are none of them in it.
\end{quote}

\begin{quote}
\textbf{Transect.} A transect answers exactly the line asked for and nothing either
side of it. A barrier two cells off the stated points is invisible to it.
\end{quote}

\begin{quote}
\textbf{Reach.} It names nothing as wrong. Scenery, a side observer island and a
shelf over the ceiling all read exactly like a shape stranded by accident, and only
the author can tell them apart.
\end{quote}

\Cref{ch:session} recounted several runs that redesigned correct boards on the
strength of a reading, and every one of them was a case of a reader taking an answer
as a verdict when the instrument had been asked a different question.

\section{The Workhorse}

\route{GET /map/\{slug\}/column} answers one or more columns from bedrock to sky
with every block named. The studio's own description of it is unambiguous about its
status:

\begin{quote}
\textsc{the workhorse}: every picture beside it is a projection, and this is what is
actually at a coordinate, which is why it is the read to reach for when a picture
and a document disagree.
\end{quote}

One run recorded that nine of its ten faults were found by reading a column and
counting courses, with the renders confirming them afterwards. When two reads
disagree, the column is the one that is not a projection.

\section{A Shape Is a Profile, Never a Point}

A column answers about a coordinate. Many claims an author wants to check are
claims about a \emph{shape}: a bank, a wall, a stair, a basin. Those need a
\textbf{transect}, which walks a polyline block by block and answers numbers.

Each station carries the ground, the storey stood on, the water and the highest
block in its column, and the step from the station before it, classed as walked,
scrambled, a barrier or a drop. Two parameters matter: \id{every} thins the
stations, and \id{beside} lists every distinct claim within a stated number of cells
of the line, which is the answer to the transect's own blind spot.

This read settles the class of question that arithmetic gets wrong. The corpus
carries several ramps whose anchors were computed in absolute terms and landed ten
or more blocks below a relief-solved hill, showing as a barrier. None was visible in
any render. Each was one transect.

\section{Four Answers in Four Units}

The \textbf{walk} is the read that prices a journey, and its design deserves note
because it refuses to summarise:

\begin{quote}
Whether the place can be reached, how far it is in blocks, how many blocks a player
must place to get there, and how many falls over the free height it takes. Four
answers in four units, none weighed against the others, so a caller reads the field
its own rule is stated in.
\end{quote}

The climb at a rise of $\delta$ costs $\delta - 1$ blocks placed, and bridging void
costs one per cell. \id{aim} picks which route is measured: the shortest
(\id{travel}), the one placing fewest blocks (\id{reach}), or the least
edge-hugging of the routes within ten blocks of the shortest (\id{comfort}).

Refusing to combine the four into a score is the right decision for a system whose
rules are stated in different units. A rule about how far a goal stands from a spawn
is stated in blocks; a rule about whether a fresh kit can reach a wool is stated in
blocks placed. A single number would serve neither.

\section{The Instruments, and What Each Is For}

\begin{figure}[htbp]
\centering
\begin{subfigure}{0.31\textwidth}
  \includegraphics[width=\linewidth]{maps/whitegape-topdown.png}
  \caption{\route{topdown}: the block on each column}
\end{subfigure}\hfill
\begin{subfigure}{0.31\textwidth}
  \includegraphics[width=\linewidth]{maps/whitegape-heightmap.png}
  \caption{\route{heightmap}: how high it stands}
\end{subfigure}\hfill
\begin{subfigure}{0.31\textwidth}
  \includegraphics[width=\linewidth]{maps/whitegape-traversability.png}
  \caption{\route{slopes}: whether it can be walked}
\end{subfigure}

\medskip
\begin{subfigure}{0.31\textwidth}
  \includegraphics[width=\linewidth]{maps/whitegape-made.png}
  \caption{\route{made}: what an author owns}
\end{subfigure}\hfill
\begin{subfigure}{0.31\textwidth}
  \includegraphics[width=\linewidth]{maps/whitegape-structure.png}
  \caption{\route{structure}: what was stamped}
\end{subfigure}\hfill
\begin{subfigure}{0.31\textwidth}
  \includegraphics[width=\linewidth]{maps/whitegape-xray.png}
  \caption{\route{xray}: what is under the surface}
\end{subfigure}
\caption{One board, six reads. Every panel is the same world asked a different
question, and no two of them could answer each other's. The point of the chapter is
in the set rather than in any one of them: an agent cannot see the map, so what it
knows is whichever of these it thought to ask for.}
\label{fig:reads-gallery}
\end{figure}


\begin{table}[htbp]
\caption{The principal reads. Every one has a \id{?format=text} twin unless noted.}
\label{tab:reads}
\small
\begin{tabularx}{\textwidth}{L{3.2cm} Y}
\toprule
\textbf{Read} & \textbf{The question it answers} \\
\midrule
\route{column} & what is actually at this coordinate, block by block \\
\route{transect} & what a line across a shape does, station by station \\
\route{slopes} & the worst step to a neighbour per cell, classed walked,
  scrambled or barrier. A cliff reads as a line, a ramp as a band crossing it, and
  an overdone relief as a page \\
\route{walk} & what a journey costs, in four units \\
\route{incline} & how steeply each cell is inclined, and how much ground stands in
  each ten degrees \\
\route{coverage} & where the ground is lived on and where it is dead, with the dead
  patches named by coordinate \\
\route{reach} & which standing ground no player can get to, and why \\
\route{themes/census} & every ground cell counted by the theme that paints it \\
\route{render/mirror} & whether a board somebody believes is symmetric actually is \\
\route{render/section} & a vertical cut with a Y scale \\
\route{render/topdown} & what was built and where, one subject per image \\
\route{plan/flow} & what the board asks of the two sides, and what that leaves
  unreached \\
\route{findings} & everything wrong with this map right now, from every gate its
  stored documents can reach \\
\bottomrule
\end{tabularx}
\end{table}

\subsection{The Incline Read, and Why It Cannot Disagree with the Paint}

\route{incline} deserves a note because of a property it has by construction. Each
cell's glyph is the \emph{tens} of its angle, so a mask cut at 20 and 45 has both
boundaries visible on the page as the step from \id{1} to \id{2} and from \id{4} to
\id{5}. Beneath the grid it prints how much ground stands in each ten degrees,
which is the reading a cut is actually chosen on.

It runs the same formula the painter resolves against, so the two cannot disagree.
The one caveat is stated in its own misleads clause: a wider \id{window} answers a
different question from the one the paint was chosen by.

\subsection{Two Reads Keep Y}

Only \route{section} and \route{column} keep elevation. Every other picture is a
projection onto the plan, which is why a top-down render can show a board that looks
complete and conceal a storey, a riser, a ramp's tread heights or a goal's
clearance. An author checking anything vertical has exactly two instruments.

\section{Findings, and the Gate That Runs Nowhere Else}

\route{GET /map/\{slug\}/findings} answers what is wrong with a map right now from
every gate its stored documents can reach, and names the gates it did \emph{not}
ask along with the route that pays for them. It is the read that answers a question
no other channel carries, because a decline published on one route's warnings is
not published on another's.

Its pair is \route{GET /map/\{slug\}/state}, which answers where the map has got
to, which artifacts it holds, and which moves those artifacts allow, each with its
route. An agent that has lost track of a map asks those two and needs nothing else.

\section{What the Reads Do Not Do}

Three limits are worth stating plainly, because each has produced a wrong report.

A read names nothing as \emph{wrong}. Several of them partition a board and hand
back the partition; deciding which part is a defect and which is scenery is the
author's. Dead-ground calibration is meaningless on a board whose objective the
coverage model does not yet understand, and on control-point boards it has
misreported badly for exactly that reason.

A read answers about the world as it stands, not about the world as played. Nothing
in this surface can say whether a map is good.

And a plan-tier read is not a world-tier read. Plan-level distances are computed
over pieces and build zones; the same journey over built ground can be shorter or
longer, and on a board with a maze in it one run measured plan-tier ratios of 2.19
to 5.31 where the built world answered 1.70 to 2.77. The plan-tier reads are a floor
and nothing more on a board authored that way.
